\documentclass[11pt]{article}
\usepackage{amsmath,amssymb,amsthm}
\usepackage[margin=1in]{geometry}
\usepackage[hidelinks]{hyperref}

\newtheorem{definition}{Definition}[section]
\newtheorem{proposition}{Proposition}[section]
\newtheorem{theorem}{Theorem}[section]
\newtheorem{lemma}{Lemma}[section]

\newcommand{\C}{\mathbb{C}}
\newcommand{\R}{\mathbb{R}}
\newcommand{\N}{\mathbb{N}}
\newcommand{\Z}{\mathbb{Z}}
\newcommand{\abs}[1]{\left|#1\right|}
\newcommand{\conj}[1]{\overline{#1}}
\newcommand{\Real}{\operatorname{Re}}
\newcommand{\Imag}{\operatorname{Im}}

\title{Complex Analysis Notes: Toward Euler's Formula}
\author{}
\date{}

\begin{document}
\maketitle
\tableofcontents

\section{Complex Numbers}

\begin{definition}[The complex numbers]
Define the complex numbers by
\[
  \boxed{
    \C=\R^{2}
    \text{ with the usual vector addition and real scalar multiplication}
  }.
\]
Equip this real vector space with multiplication
\[
  \boxed{
    (a,b)(c,d)=(ac-bd,ad+bc),
    \qquad
    1=(1,0),
    \qquad
    i=(0,1)
  }.
\]
With this multiplication,
\[
  i^{2}=(0,1)(0,1)=(-1,0).
\]
After identifying $a\in\R$ with $(a,0)\in\C$, this says $i^{2}=-1$.
Every element $(a,b)\in\R^{2}$ can be written as
\[
  (a,b)=a(1,0)+b(0,1)=a+bi.
\]
\end{definition}

\begin{definition}[Real and imaginary parts]
If $z=a+bi\in\C$, define
\[
  \Real(z):=a,
  \qquad
  \Imag(z):=b.
\]
Two complex numbers are equal precisely when their real and imaginary parts are equal:
\[
  a+bi=c+di
  \quad\Longleftrightarrow\quad
  a=c\ \text{and}\ b=d.
\]
\end{definition}

\section{Algebraic Operations}

\begin{proposition}[Complex addition in $a+bi$ notation]
For $z=a+bi$ and $w=c+di$,
\[
  z+w=(a+c)+(b+d)i.
\]
\end{proposition}

\begin{proof}
This is the usual vector addition on $\R^{2}$ written in the notation $(a,b)=a+bi$.
\end{proof}

\begin{proposition}[Complex multiplication in $a+bi$ notation]
For $z=a+bi$ and $w=c+di$,
\[
  zw=(ac-bd)+(ad+bc)i.
\]
\end{proposition}

\begin{proof}
This is exactly the multiplication
\[
  (a,b)(c,d)=(ac-bd,ad+bc)
\]
written in the notation $(a,b)=a+bi$.
\end{proof}

\begin{proposition}[Distributive law]
For all $z,w,u\in\C$,
\[
  z(w+u)=zw+zu.
\]
\end{proposition}

\begin{proof}
Let $z=a+bi$, $w=c+di$, and $u=e+fi$. Then
\[
  w+u=(c+e)+(d+f)i.
\]
Hence
\[
\begin{aligned}
  z(w+u)
    &= (a+bi)\bigl((c+e)+(d+f)i\bigr) \\
    &= \bigl(a(c+e)-b(d+f)\bigr)
       +\bigl(a(d+f)+b(c+e)\bigr)i \\
    &= (ac-bd)+(ad+bc)i+(ae-bf)+(af+be)i \\
    &= zw+zu.
\end{aligned}
\]
\end{proof}

\begin{definition}[Complex subtraction]
For $z,w\in\C$, define
\[
  z-w := z+(-w).
\]
If $z=a+bi$ and $w=c+di$, then
\[
  z-w=(a-c)+(b-d)i.
\]
\end{definition}

\section{Conjugates and Modulus}

\begin{definition}[Complex conjugate]
If $z=a+bi$, define the complex conjugate of $z$ by
\[
  \conj{z}:=a-bi.
\]
\end{definition}

\begin{proposition}[Basic conjugate identities]
For all $z,w\in\C$,
\[
  \conj{z+w}=\conj{z}+\conj{w},
  \qquad
  \conj{zw}=\conj{z}\,\conj{w},
  \qquad
  \conj{\conj{z}}=z.
\]
\end{proposition}

\begin{proof}
Each identity follows by writing $z=a+bi$ and $w=c+di$ and comparing real and imaginary parts.
\end{proof}

\begin{proposition}[Real part from conjugation]
For every $z\in\C$,
\[
  z+\conj{z}=2\Real(z).
\]
\end{proposition}

\begin{proof}
If $z=a+bi$, then
\[
  z+\conj{z}=(a+bi)+(a-bi)=2a=2\Real(z).
\]
\end{proof}

\begin{definition}[Complex modulus]
If $z=a+bi$, define the modulus of $z$ by
\[
  \abs{z}:=\sqrt{a^{2}+b^{2}}.
\]
\end{definition}

\begin{proposition}[Modulus squared]
For every $z\in\C$,
\[
  \abs{z}^{2}=z\conj{z}.
\]
\end{proposition}

\begin{proof}
If $z=a+bi$, then
\[
  z\conj{z}=(a+bi)(a-bi)=a^{2}+b^{2}=\abs{z}^{2}.
\]
\end{proof}

\begin{proposition}[Product rule for modulus]
For all $z,w\in\C$,
\[
  \abs{zw}=\abs{z}\abs{w}.
\]
\end{proposition}

\begin{proof}
Using the conjugate product identity,
\[
  \abs{zw}^{2}
  = zw\,\conj{zw}
  = zw\,\conj{z}\conj{w}
  = (z\conj{z})(w\conj{w})
  = \abs{z}^{2}\abs{w}^{2}.
\]
Both sides are nonnegative, so taking square roots gives $\abs{zw}=\abs{z}\abs{w}$.
\end{proof}

\begin{proposition}[Powers and modulus]
For every $z\in\C$ and $n\in\N$,
\[
  \abs{z^{n}}=\abs{z}^{n}.
\]
\end{proposition}

\begin{proof}
This follows by induction from the product rule:
\[
  \abs{z^{n+1}}=\abs{z^{n}z}=\abs{z^{n}}\abs{z}=\abs{z}^{n+1}.
\]
\end{proof}

\begin{proposition}[Real and imaginary parts are bounded by the modulus]
If $z\in\C$, then
\[
  \abs{\Real(z)}\le \abs{z},
  \qquad
  \abs{\Imag(z)}\le \abs{z}.
\]
\end{proposition}

\begin{proof}
If $z=a+bi$, then $a^{2}\le a^{2}+b^{2}$ and $b^{2}\le a^{2}+b^{2}$.
Taking square roots gives the result.
\end{proof}

\begin{proposition}[Triangle inequality]
For all $z,w\in\C$,
\[
  \abs{z+w}\le \abs{z}+\abs{w}.
\]
\end{proposition}

\begin{proof}
Since $\Real(\alpha)\le \abs{\alpha}$ for every $\alpha\in\C$,
\[
\begin{aligned}
  \abs{z+w}^{2}
    &= (z+w)(\conj{z}+\conj{w}) \\
    &= \abs{z}^{2}+\abs{w}^{2}+z\conj{w}+w\conj{z} \\
    &= \abs{z}^{2}+\abs{w}^{2}+2\Real(z\conj{w}) \\
    &\le \abs{z}^{2}+\abs{w}^{2}+2\abs{z\conj{w}} \\
    &= \abs{z}^{2}+\abs{w}^{2}+2\abs{z}\abs{w} \\
    &= (\abs{z}+\abs{w})^{2}.
\end{aligned}
\]
Taking square roots gives the triangle inequality.
\end{proof}

\begin{proposition}[Reverse triangle inequality]
For all $z,w\in\C$,
\[
  \bigl|\abs{z}-\abs{w}\bigr|\le \abs{z-w}.
\]
\end{proposition}

\begin{proof}
By the triangle inequality,
\[
  \abs{z}=\abs{(z-w)+w}\le \abs{z-w}+\abs{w},
\]
so $\abs{z}-\abs{w}\le \abs{z-w}$.
Interchanging $z$ and $w$ gives
\[
  \abs{w}-\abs{z}\le \abs{w-z}=\abs{z-w}.
\]
Therefore $\bigl|\abs{z}-\abs{w}\bigr|\le \abs{z-w}$.
\end{proof}

\section{Reciprocals and Division}

\begin{definition}[Complex reciprocal]
If $z\in\C$ and $z\ne0$, define
\[
  \frac{1}{z}:=\frac{\conj{z}}{\abs{z}^{2}}.
\]
\end{definition}

\begin{proposition}[The reciprocal is the multiplicative inverse]
If $z\ne0$, then
\[
  z\cdot\frac{1}{z}=1.
\]
\end{proposition}

\begin{proof}
Using $\abs{z}^{2}=z\conj{z}$,
\[
  z\cdot\frac{1}{z}
  = z\cdot\frac{\conj{z}}{\abs{z}^{2}}
  = \frac{z\conj{z}}{\abs{z}^{2}}
  = 1.
\]
\end{proof}

\begin{proposition}[$\C$ is a field]
With the addition and multiplication defined above, $\C$ is a field.
\end{proposition}

\begin{proof}
The additive identity is $0=0+0i$, and the multiplicative identity is $1=1+0i$.
Associativity and commutativity follow by direct expansion from the corresponding real-number laws, and distributivity was proved above.
If $z=a+bi$, then the additive inverse is
\[
  -z=(-a)+(-b)i.
\]
If $z\ne0$, then the reciprocal theorem gives a multiplicative inverse.
Therefore $\C$ is a field.
\end{proof}

\begin{definition}[Complex division]
For $w,z\in\C$ with $z\ne0$, define
\[
  \frac{w}{z}:=w\cdot\frac{1}{z}
  =\frac{w\conj{z}}{\abs{z}^{2}}.
\]
\end{definition}

\section{Complex Sequences}

\begin{definition}[Complex sequence]
A complex sequence is a function
\[
  z:\N\to\C.
\]
We write $z_{n}$ for $z(n)$.
If $z_{n}=a_{n}+ib_{n}$, then $(a_{n})$ and $(b_{n})$ are real sequences, and
\[
  a_{n}=\Real(z_{n}),
  \qquad
  b_{n}=\Imag(z_{n}).
\]
\end{definition}

\begin{definition}[$\varepsilon$-$N$ convergence in $\C$]
Let $(z_{n})$ be a complex sequence and let $z\in\C$.
We say
\[
  z_{n}\to z
\]
if for every $\varepsilon>0$, there exists $N\in\N$ such that
\[
  n>N \quad\Longrightarrow\quad \abs{z_{n}-z}<\varepsilon.
\]
\end{definition}

\begin{proposition}[Convergence as distance to the limit]
Let $(z_{n})$ be a complex sequence and let $z\in\C$.
Then
\[
  z_{n}\to z
  \quad\Longleftrightarrow\quad
  \abs{z_{n}-z}\to0
\]
as a real sequence.
\end{proposition}

\begin{proof}
The sequence $\bigl(\abs{z_{n}-z}\bigr)$ is a real sequence because each term is a nonnegative real number.

First suppose $z_{n}\to z$ in $\C$.
By definition, for every $\varepsilon>0$, there exists $N\in\N$ such that
\[
  n>N\quad\Longrightarrow\quad \abs{z_{n}-z}<\varepsilon.
\]
This is exactly the statement that the real sequence $\abs{z_{n}-z}$ converges to $0$.

Conversely, suppose $\abs{z_{n}-z}\to0$ as a real sequence.
Then for every $\varepsilon>0$, there exists $N\in\N$ such that
\[
  n>N\quad\Longrightarrow\quad
  \bigl|\abs{z_{n}-z}-0\bigr|<\varepsilon.
\]
Since $\abs{z_{n}-z}\ge0$,
\[
  \bigl|\abs{z_{n}-z}-0\bigr|=\abs{z_{n}-z}.
\]
Thus
\[
  n>N\quad\Longrightarrow\quad \abs{z_{n}-z}<\varepsilon,
\]
which is precisely $z_{n}\to z$ in $\C$.
\end{proof}

\begin{theorem}[Componentwise convergence]
Let $z_{n}=a_{n}+ib_{n}$ and $z=a+ib$.
Then
\[
  z_{n}\to z
  \quad\Longleftrightarrow\quad
  a_{n}\to a\ \text{and}\ b_{n}\to b.
\]
Equivalently,
\[
  z_{n}\to z
  \quad\Longleftrightarrow\quad
  \Real(z_{n})\to\Real(z)\ \text{and}\ \Imag(z_{n})\to\Imag(z).
\]
\end{theorem}

\begin{proof}
First suppose $z_{n}\to z$.
Since
\[
  \abs{a_{n}-a}=\abs{\Real(z_{n}-z)}\le \abs{z_{n}-z},
  \qquad
  \abs{b_{n}-b}=\abs{\Imag(z_{n}-z)}\le \abs{z_{n}-z},
\]
the $\varepsilon$-$N$ definition for $z_{n}\to z$ implies $a_{n}\to a$ and $b_{n}\to b$.

Conversely, suppose $a_{n}\to a$ and $b_{n}\to b$.
Let $\varepsilon>0$.
Choose $N_{1}$ such that
\[
  n>N_{1}\quad\Longrightarrow\quad \abs{a_{n}-a}<\frac{\varepsilon}{2},
\]
and choose $N_{2}$ such that
\[
  n>N_{2}\quad\Longrightarrow\quad \abs{b_{n}-b}<\frac{\varepsilon}{2}.
\]
Let $N=\max\{N_{1},N_{2}\}$.
Then for $n>N$,
\[
\begin{aligned}
  \abs{z_{n}-z}
    &=\abs{(a_{n}-a)+i(b_{n}-b)} \\
    &=\sqrt{(a_{n}-a)^{2}+(b_{n}-b)^{2}} \\
    &\le \abs{a_{n}-a}+\abs{b_{n}-b} \\
    &<\varepsilon.
\end{aligned}
\]
Therefore $z_{n}\to z$.
\end{proof}

\begin{theorem}[Uniqueness of limits]
If $z_{n}\to z$ and $z_{n}\to w$, then $z=w$.
\end{theorem}

\begin{proof}
Write
\[
  z=a+ib,
  \qquad
  w=c+id,
  \qquad
  z_{n}=a_{n}+ib_{n}.
\]
By componentwise convergence,
\[
  a_{n}\to a,\quad b_{n}\to b,
  \qquad
  a_{n}\to c,\quad b_{n}\to d.
\]
Real limits are unique, so $a=c$ and $b=d$.
Therefore $z=w$.
\end{proof}

\begin{theorem}[Convergent sequences are bounded]
If $z_{n}\to z$, then $(z_{n})$ is bounded: there exists $M>0$ such that
\[
  \abs{z_{n}}\le M
  \qquad
  \forall n\in\N.
\]
\end{theorem}

\begin{proof}
Since $z_{n}\to z$, choose $N\in\N$ such that
\[
  n>N\quad\Longrightarrow\quad \abs{z_{n}-z}<1.
\]
For $n>N$, the triangle inequality gives
\[
  \abs{z_{n}}
  =\abs{(z_{n}-z)+z}
  \le \abs{z_{n}-z}+\abs{z}
  <1+\abs{z}.
\]
Now define
\[
  M:=\max\{\abs{z_{1}},\abs{z_{2}},\dots,\abs{z_{N}},1+\abs{z}\}.
\]
Then $\abs{z_{n}}\le M$ for every $n\in\N$.
\end{proof}

\begin{definition}[Subsequence]
Let $(z_{n})$ be a complex sequence.
A subsequence is a sequence of the form
\[
  z_{n_{k}},
\]
where
\[
  n_{1}<n_{2}<n_{3}<\cdots
\]
is a strictly increasing sequence of natural numbers.
\end{definition}

\begin{proposition}[Subsequences preserve limits]
If $z_{n}\to z$, then every subsequence $z_{n_{k}}$ also converges to $z$.
\end{proposition}

\begin{proof}
Let $\varepsilon>0$.
Choose $N\in\N$ such that
\[
  n>N\quad\Longrightarrow\quad \abs{z_{n}-z}<\varepsilon.
\]
Since $n_{k}\ge k$, every $k>N$ satisfies $n_{k}>N$.
Thus
\[
  k>N\quad\Longrightarrow\quad \abs{z_{n_{k}}-z}<\varepsilon.
\]
Therefore $z_{n_{k}}\to z$.
\end{proof}

\begin{definition}[Cauchy sequence in $\C$]
A complex sequence $(z_{n})$ is Cauchy if for every $\varepsilon>0$, there exists $N\in\N$ such that
\[
  m,n>N\quad\Longrightarrow\quad \abs{z_{n}-z_{m}}<\varepsilon.
\]
\end{definition}

\begin{theorem}[Cauchy criterion in $\C$]
A complex sequence $(z_{n})$ converges if and only if it is Cauchy.
\end{theorem}

\begin{proof}
Write $z_{n}=a_{n}+ib_{n}$.

Suppose $z_{n}\to z$.
Every convergent sequence is Cauchy: given $\varepsilon>0$, choose $N$ such that
\[
  n>N\quad\Longrightarrow\quad \abs{z_{n}-z}<\frac{\varepsilon}{2}.
\]
Then for $m,n>N$,
\[
  \abs{z_{n}-z_{m}}
  \le \abs{z_{n}-z}+\abs{z_{m}-z}
  <\varepsilon.
\]

Conversely, suppose $(z_{n})$ is Cauchy.
Since
\[
  \abs{a_{n}-a_{m}}\le \abs{z_{n}-z_{m}},
  \qquad
  \abs{b_{n}-b_{m}}\le \abs{z_{n}-z_{m}},
\]
the real sequences $(a_{n})$ and $(b_{n})$ are Cauchy.
Because $\R$ is complete, there exist $a,b\in\R$ such that
\[
  a_{n}\to a,
  \qquad
  b_{n}\to b.
\]
By componentwise convergence, $z_{n}\to a+ib$.
\end{proof}

\begin{theorem}[Limit algebra]
Suppose $z_{n}\to z$ and $w_{n}\to w$ in $\C$, and let $\alpha\in\C$.
Then
\[
  z_{n}+w_{n}\to z+w,
  \qquad
  z_{n}-w_{n}\to z-w,
  \qquad
  \alpha z_{n}\to \alpha z,
  \qquad
  z_{n}w_{n}\to zw.
\]
If $w\ne0$ and $w_{n}\ne0$ eventually, then
\[
  \frac{z_{n}}{w_{n}}\to\frac{z}{w}.
\]
\end{theorem}

\begin{proof}
For sums,
\[
  \abs{(z_{n}+w_{n})-(z+w)}
  \le \abs{z_{n}-z}+\abs{w_{n}-w}\to0.
\]
Subtraction follows from addition and scalar multiplication by $-1$.
For scalar multiplication,
\[
  \abs{\alpha z_{n}-\alpha z}
  =\abs{\alpha}\abs{z_{n}-z}\to0.
\]
For products,
\[
\begin{aligned}
  z_{n}w_{n}-zw
    &=z_{n}w_{n}-z_{n}w+z_{n}w-zw \\
    &=z_{n}(w_{n}-w)+(z_{n}-z)w.
\end{aligned}
\]
Since $(z_{n})$ is convergent, it is bounded; choose $M>0$ such that $\abs{z_{n}}\le M$ for all $n$.
Then
\[
  \abs{z_{n}w_{n}-zw}
  \le M\abs{w_{n}-w}+\abs{w}\abs{z_{n}-z}\to0.
\]

For reciprocals, assume $w\ne0$.
Since $w_{n}\to w$, choose $N$ such that
\[
  n>N\quad\Longrightarrow\quad \abs{w_{n}-w}<\frac{\abs{w}}{2}.
\]
Then for $n>N$,
\[
  \abs{w_{n}}\ge \abs{w}-\abs{w_{n}-w}>\frac{\abs{w}}{2}.
\]
Therefore
\[
  \abs{\frac{1}{w_{n}}-\frac{1}{w}}
  =\abs{\frac{w-w_{n}}{w_{n}w}}
  =\frac{\abs{w_{n}-w}}{\abs{w_{n}}\abs{w}}
  \le \frac{2}{\abs{w}^{2}}\abs{w_{n}-w}\to0.
\]
Thus $1/w_{n}\to1/w$, and the quotient rule follows from the product rule:
\[
  \frac{z_{n}}{w_{n}}=z_{n}\cdot\frac{1}{w_{n}}\to z\cdot\frac{1}{w}.
\]
\end{proof}

\begin{proposition}[Limit of the modulus]
If $z_{n}\to z$, then
\[
  \abs{z_{n}}\to \abs{z}.
\]
\end{proposition}

\begin{proof}
By the reverse triangle inequality,
\[
  \bigl|\abs{z_{n}}-\abs{z}\bigr|\le \abs{z_{n}-z}.
\]
Since $\abs{z_{n}-z}\to0$, it follows that $\abs{z_{n}}\to\abs{z}$.
\end{proof}

\begin{proposition}[Limit of the conjugate]
If $z_{n}\to z$, then
\[
  \conj{z_{n}}\to \conj{z}.
\]
\end{proposition}

\begin{proof}
Since
\[
  \abs{\conj{z_{n}}-\conj{z}}
  =\abs{\conj{z_{n}-z}}
  =\abs{z_{n}-z},
\]
we have $\conj{z_{n}}\to\conj{z}$.
\end{proof}

\section{Complex Series}

\begin{definition}[Partial sum sequence]
Let $(z_{n})_{n=0}^{\infty}$ be a complex sequence.
The associated partial sum sequence $(s_{N})_{N=0}^{\infty}$ is defined by
\[
  s_{N}:=\sum_{n=0}^{N} z_{n}
  =z_{0}+z_{1}+\cdots+z_{N}.
\]
\end{definition}

\begin{definition}[Infinite sum]
Let $(z_{n})_{n=0}^{\infty}$ be a complex sequence with partial sums
\[
  s_{N}=\sum_{n=0}^{N}z_{n}.
\]
If the partial sum sequence converges to some $S\in\C$, then the infinite series converges and we write
\[
  \sum_{n=0}^{\infty}z_{n}=S.
\]
Equivalently,
\[
  \sum_{n=0}^{\infty}z_{n}=S
  \quad\Longleftrightarrow\quad
  \lim_{N\to\infty}\sum_{n=0}^{N}z_{n}=S.
\]
\end{definition}

\begin{definition}[Absolute convergence]
A complex series $\sum_{n=0}^{\infty}z_{n}$ is absolutely convergent if the real series
\[
  \sum_{n=0}^{\infty}\abs{z_{n}}
\]
converges.
\end{definition}

\begin{theorem}[Absolute convergence implies convergence]
If $\sum_{n=0}^{\infty}z_{n}$ is absolutely convergent, then $\sum_{n=0}^{\infty}z_{n}$ converges.
\end{theorem}

\begin{proof}
Let
\[
  s_{N}:=\sum_{n=0}^{N}z_{n},
  \qquad
  t_{N}:=\sum_{n=0}^{N}\abs{z_{n}}.
\]
Since $\sum\abs{z_{n}}$ converges in $\R$, the real partial sums $(t_{N})$ are Cauchy.
Let $\varepsilon>0$.
Choose $N\in\N$ such that
\[
  q>p>N
  \quad\Longrightarrow\quad
  \sum_{n=p+1}^{q}\abs{z_{n}}
  =t_{q}-t_{p}<\varepsilon.
\]
Then for $q>p>N$,
\[
  \abs{s_{q}-s_{p}}
  =\abs{\sum_{n=p+1}^{q}z_{n}}
  \le \sum_{n=p+1}^{q}\abs{z_{n}}
  <\varepsilon.
\]
Thus $(s_{N})$ is Cauchy in $\C$, so it converges by the Cauchy criterion in $\C$.
Therefore $\sum z_{n}$ converges.
\end{proof}

\begin{theorem}[Reordering absolutely convergent series]
Suppose $\sum_{n=0}^{\infty}z_{n}$ is absolutely convergent and
\[
  \sum_{n=0}^{\infty}z_{n}=S.
\]
If $\sigma:\N_{0}\to\N_{0}$ is a bijection, then
\[
  \sum_{n=0}^{\infty}z_{\sigma(n)}
\]
also converges absolutely, and
\[
  \sum_{n=0}^{\infty}z_{\sigma(n)}=S.
\]
\end{theorem}

\begin{proof}
First, the reordered absolute series converges.
For every $N$, the finite set $\{\sigma(0),\dots,\sigma(N)\}$ is contained in $\{0,\dots,M\}$ for some $M$, so
\[
  \sum_{n=0}^{N}\abs{z_{\sigma(n)}}
  \le \sum_{n=0}^{M}\abs{z_{n}}
  \le \sum_{n=0}^{\infty}\abs{z_{n}}.
\]
Thus the partial sums of $\sum\abs{z_{\sigma(n)}}$ are increasing and bounded, hence converge.

Now let
\[
  s_{N}:=\sum_{n=0}^{N}z_{n},
  \qquad
  r_{N}:=\sum_{n=0}^{N}z_{\sigma(n)}.
\]
Let $\varepsilon>0$.
Since $\sum\abs{z_{n}}$ converges and $s_{K}\to S$, choose $K$ such that
\[
  \sum_{j=K+1}^{\infty}\abs{z_{j}}<\varepsilon
  \qquad\text{and}\qquad
  \abs{s_{K}-S}<\varepsilon.
\]
Choose $N$ large enough that
\[
  \{0,1,\dots,K\}\subseteq \{\sigma(0),\sigma(1),\dots,\sigma(N)\}.
\]
For $m>N$, the difference $r_{m}-s_{K}$ is a finite sum of terms $z_{j}$ with $j>K$.
Therefore
\[
  \abs{r_{m}-s_{K}}
  \le \sum_{\substack{j>K\\ j\in\{\sigma(0),\dots,\sigma(m)\}}}\abs{z_{j}}
  \le \sum_{j=K+1}^{\infty}\abs{z_{j}}
  \le \varepsilon.
\]
Then for $m>N$,
\[
  \abs{r_{m}-S}
  \le \abs{r_{m}-s_{K}}+\abs{s_{K}-S}
  <2\varepsilon.
\]
Hence $r_{m}\to S$, so the reordered series converges to the same sum.
\end{proof}

\begin{theorem}[Ratio test]
Let $\sum_{n=0}^{\infty}z_{n}$ be a complex series, and suppose $z_{n}\ne0$ for all sufficiently large $n$.
If
\[
  \lim_{n\to\infty}\frac{\abs{z_{n+1}}}{\abs{z_{n}}}=L,
\]
then:
\[
\begin{array}{ll}
  L<1 & \Longrightarrow \sum_{n=0}^{\infty}z_{n}\ \text{converges absolutely},\\[4pt]
  L>1 & \Longrightarrow \sum_{n=0}^{\infty}z_{n}\ \text{diverges},\\[4pt]
  L=1 & \Longrightarrow \text{the test is inconclusive}.
\end{array}
\]
\end{theorem}

\begin{proof}
First suppose $L<1$.
Choose $r$ such that
\[
  L<r<1.
\]
Since $\abs{z_{n+1}}/\abs{z_{n}}\to L$, there exists $N\in\N$ such that
\[
  n>N
  \quad\Longrightarrow\quad
  \frac{\abs{z_{n+1}}}{\abs{z_{n}}}<r.
\]
Thus, for $k\ge1$,
\[
  \abs{z_{N+k}}
  < r^{k}\abs{z_{N}}.
\]
Therefore
\[
  \sum_{k=1}^{\infty}\abs{z_{N+k}}
  \le \abs{z_{N}}\sum_{k=1}^{\infty}r^{k}<\infty.
\]
Adding finitely many initial terms does not affect convergence, so $\sum\abs{z_{n}}$ converges.
Hence $\sum z_{n}$ converges absolutely.

Now suppose $L>1$.
Choose $r>1$ such that
\[
  1<r<L.
\]
Then for all sufficiently large $n$,
\[
  \frac{\abs{z_{n+1}}}{\abs{z_{n}}}>r,
\]
so $\abs{z_{n+1}}>r\abs{z_{n}}$ eventually.
In particular, $\abs{z_{n}}$ does not converge to $0$.
Hence $z_{n}$ does not converge to $0$, so the series $\sum z_{n}$ diverges.

When $L=1$, the test gives no conclusion.
For example, $\sum 1/n$ diverges but $\sum 1/n^{2}$ converges, and both have ratio limit $1$.
\end{proof}

\section{The Complex Exponential and Euler's Formula}

\begin{theorem}[Binomial theorem over $\C$]
For all $z,w\in\C$ and $n\in\N$,
\[
  (z+w)^{n}=\sum_{k=0}^{n}\binom{n}{k}z^{k}w^{n-k}.
\]
\end{theorem}

\begin{proof}
The proof is the same induction as over $\R$, using that $\C$ is a commutative field.
For $n=1$, the statement is immediate.
Assume it holds for $n$.
Then
\[
\begin{aligned}
  (z+w)^{n+1}
    &= (z+w)\sum_{k=0}^{n}\binom{n}{k}z^{k}w^{n-k} \\
    &= \sum_{k=0}^{n}\binom{n}{k}z^{k+1}w^{n-k}
       +\sum_{k=0}^{n}\binom{n}{k}z^{k}w^{n+1-k}.
\end{aligned}
\]
Collecting equal powers and using Pascal's identity gives
\[
  (z+w)^{n+1}
  =\sum_{k=0}^{n+1}\binom{n+1}{k}z^{k}w^{n+1-k}.
\]
\end{proof}

\begin{theorem}[Cauchy product theorem]
Suppose $\sum_{n=0}^{\infty}a_{n}$ and $\sum_{n=0}^{\infty}b_{n}$ are absolutely convergent complex series, and define
\[
  c_{n}:=\sum_{k=0}^{n}a_{k}b_{n-k}.
\]
Then $\sum c_{n}$ converges absolutely and
\[
  \sum_{n=0}^{\infty}c_{n}
  =
  \left(\sum_{n=0}^{\infty}a_{n}\right)
  \left(\sum_{n=0}^{\infty}b_{n}\right).
\]
\end{theorem}

\begin{proof}
Let
\[
  A_{*}:=\sum_{n=0}^{\infty}\abs{a_{n}},
  \qquad
  B_{*}:=\sum_{n=0}^{\infty}\abs{b_{n}}.
\]
For every $N$,
\[
  \sum_{n=0}^{N}\abs{c_{n}}
  \le
  \sum_{n=0}^{N}\sum_{k=0}^{n}\abs{a_{k}}\abs{b_{n-k}}
  \le A_{*}B_{*}.
\]
Thus $\sum c_{n}$ is absolutely convergent.

Let
\[
  A=\sum_{n=0}^{\infty}a_{n},
  \qquad
  B=\sum_{n=0}^{\infty}b_{n},
\]
and define
\[
  A_{N}:=\sum_{n=0}^{N}a_{n},
  \qquad
  B_{N}:=\sum_{n=0}^{N}b_{n},
  \qquad
  C_{N}:=\sum_{n=0}^{N}c_{n}.
\]
Then
\[
  A_{N}B_{N}
  =
  \sum_{\substack{0\le k\le N\\0\le j\le N}}a_{k}b_{j},
  \qquad
  C_{N}
  =
  \sum_{\substack{k,j\ge0\\k+j\le N}}a_{k}b_{j}.
\]
Choose $K$ so large that both tails
\[
  \sum_{k>K}\abs{a_{k}},
  \qquad
  \sum_{j>K}\abs{b_{j}}
\]
are small. If $N>2K$, then every pair satisfying $0\le k,j\le N$ and $k+j>N$ has either $k>K$ or $j>K$.
Hence
\[
  \abs{A_{N}B_{N}-C_{N}}
  \le
  B_{*}\sum_{k>K}\abs{a_{k}}
  +
  A_{*}\sum_{j>K}\abs{b_{j}},
\]
which can be made arbitrarily small.
Since $A_{N}\to A$ and $B_{N}\to B$, we have $A_{N}B_{N}\to AB$.
Therefore $C_{N}\to AB$.
\end{proof}

\begin{definition}[The complex exponential]
For $z\in\C$, define
\[
  \exp(z):=\sum_{n=0}^{\infty}\frac{z^{n}}{n!}.
\]
\end{definition}

\begin{proposition}[The exponential series converges absolutely]
For every $z\in\C$, the series
\[
  \sum_{n=0}^{\infty}\frac{z^{n}}{n!}
\]
converges absolutely.
\end{proposition}

\begin{proof}
If $z=0$, then
\[
  \sum_{n=0}^{\infty}\frac{z^{n}}{n!}=1+0+0+\cdots
\]
is absolutely convergent.

Now suppose $z\ne0$.
Let
\[
  a_{n}:=\frac{z^{n}}{n!}.
\]
Then $a_{n}\ne0$ for every $n$, and
\[
  \frac{\abs{a_{n+1}}}{\abs{a_{n}}}
  =
  \frac{\abs{z^{n+1}}}{(n+1)!}\frac{n!}{\abs{z^{n}}}
  =
  \frac{\abs{z}}{n+1}
  \to0.
\]
By the ratio test, $\sum a_{n}$ converges absolutely.
\end{proof}

\begin{theorem}[Exponential addition law]
For all $z,w\in\C$,
\[
  \exp(z+w)=\exp(z)\exp(w).
\]
\end{theorem}

\begin{proof}
The exponential series for $z$ and $w$ converge absolutely, so their product is given by the Cauchy product:
\[
\begin{aligned}
  \exp(z)\exp(w)
    &=
    \left(\sum_{n=0}^{\infty}\frac{z^{n}}{n!}\right)
    \left(\sum_{n=0}^{\infty}\frac{w^{n}}{n!}\right) \\
    &=
    \sum_{n=0}^{\infty}
    \sum_{k=0}^{n}
    \frac{z^{k}}{k!}\frac{w^{n-k}}{(n-k)!}.
\end{aligned}
\]
Since
\[
  \frac{1}{k!(n-k)!}
  =
  \frac{1}{n!}\binom{n}{k},
\]
we get
\[
\begin{aligned}
  \exp(z)\exp(w)
    &=
    \sum_{n=0}^{\infty}\frac{1}{n!}
    \sum_{k=0}^{n}\binom{n}{k}z^{k}w^{n-k} \\
    &=
    \sum_{n=0}^{\infty}\frac{(z+w)^{n}}{n!} \\
    &= \exp(z+w),
\end{aligned}
\]
where the binomial theorem was used in the second-to-last line.
\end{proof}

\begin{definition}[Sine and cosine by power series]
For $z\in\C$, define
\[
  \cos z
  :=\sum_{n=0}^{\infty}(-1)^{n}\frac{z^{2n}}{(2n)!},
  \qquad
  \sin z
  :=\sum_{n=0}^{\infty}(-1)^{n}\frac{z^{2n+1}}{(2n+1)!}.
\]
\end{definition}

\begin{proposition}[The sine and cosine series converge absolutely]
For every $z\in\C$, the series defining $\cos z$ and $\sin z$ converge absolutely.
\end{proposition}

\begin{proof}
If $z=0$, convergence is immediate.
Assume $z\ne0$.
For the cosine series, let
\[
  a_{n}:=(-1)^{n}\frac{z^{2n}}{(2n)!}.
\]
Then
\[
  \frac{\abs{a_{n+1}}}{\abs{a_{n}}}
  =
  \frac{\abs{z}^{2}}{(2n+2)(2n+1)}
  \to0.
\]
For the sine series, let
\[
  b_{n}:=(-1)^{n}\frac{z^{2n+1}}{(2n+1)!}.
\]
Then
\[
  \frac{\abs{b_{n+1}}}{\abs{b_{n}}}
  =
  \frac{\abs{z}^{2}}{(2n+3)(2n+2)}
  \to0.
\]
Both series converge absolutely by the ratio test.
\end{proof}

\begin{theorem}[Euler's formula over $\C$]
For every $z\in\C$,
\[
  \exp(iz)=\cos z+i\sin z.
\]
\end{theorem}

\begin{proof}
By definition,
\[
  \exp(iz)=\sum_{n=0}^{\infty}\frac{(iz)^{n}}{n!}.
\]
This exponential series is absolutely convergent, so its terms may be reordered and grouped without changing the sum.
Separate the even and odd powers:
\[
\begin{aligned}
  \exp(iz)
    &= \sum_{n=0}^{\infty}\frac{(iz)^{2n}}{(2n)!}
       +\sum_{n=0}^{\infty}\frac{(iz)^{2n+1}}{(2n+1)!}.
\end{aligned}
\]
Since
\[
  i^{2n}=(-1)^{n},
  \qquad
  i^{2n+1}=i(-1)^{n},
\]
we get
\[
\begin{aligned}
  \exp(iz)
    &= \sum_{n=0}^{\infty}(-1)^{n}\frac{z^{2n}}{(2n)!}
       + i\sum_{n=0}^{\infty}(-1)^{n}\frac{z^{2n+1}}{(2n+1)!} \\
    &= \cos z+i\sin z.
\end{aligned}
\]
\end{proof}

\begin{theorem}[Exponential formulas for sine and cosine]
For every $z\in\C$,
\[
  \boxed{\cos z=\frac{\exp(iz)+\exp(-iz)}{2}},
  \qquad
  \boxed{\sin z=\frac{\exp(iz)-\exp(-iz)}{2i}}.
\]
\end{theorem}

\begin{proof}
From the power-series definitions,
\[
  \cos(-z)=\cos z,
  \qquad
  \sin(-z)=-\sin z.
\]
Euler's formula applied to $z$ and $-z$ gives
\[
  \exp(iz)=\cos z+i\sin z,
  \qquad
  \exp(-iz)=\cos z-i\sin z.
\]
Adding these identities gives
\[
  \exp(iz)+\exp(-iz)=2\cos z.
\]
Subtracting gives
\[
  \exp(iz)-\exp(-iz)=2i\sin z.
\]
\end{proof}

\begin{theorem}[Euler's identity for real angles]
For every $\theta\in\R$,
\[
  \exp(i\theta)=\cos\theta+i\sin\theta.
\]
\end{theorem}

\begin{proof}
This is the special case of Euler's formula over $\C$ with $z=\theta\in\R$.
\end{proof}

\begin{theorem}[Real and imaginary parts of $\exp z$]
Let $z=a+ib$ with $a,b\in\R$.
Then
\[
  \exp(z)=\exp(a)(\cos b+i\sin b).
\]
Consequently,
\[
  \Real(\exp z)=\exp(a)\cos b,
  \qquad
  \Imag(\exp z)=\exp(a)\sin b.
\]
Equivalently,
\[
  \Real(\exp z)=\exp(\Real z)\cos(\Imag z),
  \qquad
  \Imag(\exp z)=\exp(\Real z)\sin(\Imag z).
\]
\end{theorem}

\begin{proof}
Using the exponential addition law and Euler's identity for real angles,
\[
  \exp(z)=\exp(a+ib)=\exp(a)\exp(ib)=\exp(a)(\cos b+i\sin b).
\]
Since $a,b\in\R$, the power-series definitions give $\exp(a),\cos b,\sin b\in\R$, so
\[
  \exp(z)=\exp(a)\cos b+i\exp(a)\sin b.
\]
Reading off real and imaginary parts gives the result.
\end{proof}

\begin{theorem}[Modulus of the complex exponential]
For every $z\in\C$,
\[
  \abs{\exp(z)}=\exp(\Real z).
\]
In particular, for every $\theta\in\R$,
\[
  \abs{\exp(i\theta)}=1.
\]
\end{theorem}

\begin{proof}
Write $z=a+ib$ with $a,b\in\R$.
From the previous theorem,
\[
  \exp(z)=\exp(a)(\cos b+i\sin b).
\]
Thus
\[
\begin{aligned}
  \abs{\exp(z)}
    &= \abs{\exp(a)}\abs{\cos b+i\sin b} \\
    &= \exp(a)\sqrt{\cos^{2}b+\sin^{2}b} \\
    &= \exp(a).
\end{aligned}
\]
Here we use the standard real-analysis facts that $\exp(a)>0$ for real $a$ and $\cos^{2}b+\sin^{2}b=1$.
Since $a=\Real z$, this proves $\abs{\exp(z)}=\exp(\Real z)$.
If $z=i\theta$, then $\Real(i\theta)=0$, so
\[
  \abs{\exp(i\theta)}=\exp(0)=1.
\]
\end{proof}

\section{Argument, Polar Form, and De Moivre's Theorem}

\begin{definition}[Argument]
Let $z\in\C$ with $z\ne0$.
An argument of $z$ is any real number $\theta$ such that
\[
  z=\abs{z}(\cos\theta+i\sin\theta).
\]
The set of all arguments of $z$ is denoted
\[
  \arg(z).
\]
Arguments are not unique: if $\theta\in\arg(z)$, then
\[
  \arg(z)=\{\theta+2\pi k:k\in\Z\}.
\]
\end{definition}

\begin{definition}[Polar form]
If $z\ne0$ and $\theta\in\arg(z)$, then the polar form of $z$ is
\[
  z=r(\cos\theta+i\sin\theta),
  \qquad
  r=\abs{z}>0.
\]
Using Euler's formula, this can also be written as
\[
  z=re^{i\theta}.
\]
\end{definition}

\begin{proposition}[Existence of polar form]
Every nonzero complex number has a polar form.
\end{proposition}

\begin{proof}
Let $z=a+ib\ne0$ and set $r=\abs{z}=\sqrt{a^{2}+b^{2}}$.
Then $r>0$, and
\[
  \left(\frac{a}{r}\right)^{2}+\left(\frac{b}{r}\right)^{2}=1.
\]
By the standard real parametrization of the unit circle, there exists $\theta\in\R$ such that
\[
  \cos\theta=\frac{a}{r},
  \qquad
  \sin\theta=\frac{b}{r}.
\]
Therefore
\[
  z=a+ib=r(\cos\theta+i\sin\theta)=re^{i\theta}.
\]
\end{proof}

\begin{proposition}[Multiplication in polar form]
If
\[
  z=re^{i\theta},
  \qquad
  w=se^{i\phi},
\]
with $r,s>0$, then
\[
  zw=rs\,e^{i(\theta+\phi)}.
\]
\end{proposition}

\begin{proof}
Using the exponential addition law,
\[
  zw=(re^{i\theta})(se^{i\phi})=rs\,e^{i\theta}e^{i\phi}=rs\,e^{i(\theta+\phi)}.
\]
\end{proof}

\begin{theorem}[De Moivre's theorem]
For every $\theta\in\R$ and $n\in\N$,
\[
  (\cos\theta+i\sin\theta)^{n}
  =
  \cos(n\theta)+i\sin(n\theta).
\]
Equivalently,
\[
  (e^{i\theta})^{n}=e^{in\theta}.
\]
\end{theorem}

\begin{proof}
By Euler's formula,
\[
  \cos\theta+i\sin\theta=e^{i\theta}.
\]
Using the exponential addition law repeatedly,
\[
  (e^{i\theta})^{n}
  =
  \underbrace{e^{i\theta}\cdots e^{i\theta}}_{n\ \text{factors}}
  =
  e^{i(\theta+\cdots+\theta)}
  =
  e^{in\theta}.
\]
Applying Euler's formula again gives
\[
  e^{in\theta}=\cos(n\theta)+i\sin(n\theta).
\]
\end{proof}

\newpage
\section*{Common Forms and Functions}

Let
\[
  z=a+bi=re^{i\theta}=r(\cos\theta+i\sin\theta),
  \qquad
  r=\abs{z}>0,
  \qquad
  \theta\in\arg(z).
\]

\[
\begin{array}{c|c|c}
\text{Object} & \text{Cartesian form} & \text{Polar/exponential form} \\
\hline
z & a+bi & re^{i\theta} \\
\abs{z} & \sqrt{a^{2}+b^{2}} & r \\
\arg(z) & \theta\ \text{s.t. } \cos\theta=a/r,\ \sin\theta=b/r
        & \{\theta+2\pi k:k\in\Z\} \\
\dfrac{1}{z} & \dfrac{a-bi}{a^{2}+b^{2}} & r^{-1}e^{-i\theta} \\
\exp(z) & \exp(a)(\cos b+i\sin b) & \exp(a)e^{ib} \\
z^{n} & (a+bi)^{n} & r^{n}e^{in\theta} \\
\cos z & \dfrac{\exp(iz)+\exp(-iz)}{2} & \dfrac{e^{iz}+e^{-iz}}{2} \\
\sin z & \dfrac{\exp(iz)-\exp(-iz)}{2i} & \dfrac{e^{iz}-e^{-iz}}{2i}
\end{array}
\]

\[
\begin{array}{c|c|c}
\text{Expression} & \Real(\cdot) & \Imag(\cdot) \\
\hline
z=a+bi & a & b \\
\conj{z}=a-bi & a & -b \\
\exp(a+bi) & \exp(a)\cos b & \exp(a)\sin b \\
r e^{i\theta} & r\cos\theta & r\sin\theta \\
z^{n}=r^{n}e^{in\theta} & r^{n}\cos(n\theta) & r^{n}\sin(n\theta)
\end{array}
\]


\end{document}
